\subsection{Lý do chọn đề tài}
Sự phát triển mạnh mẽ của thương mại điện tử trong những năm gần đây đã làm thay đổi toàn diện cách thức mua sắm và tiêu dùng của người dùng. Đặc biệt trong lĩnh vực bán lẻ điện thoại và thiết bị công nghệ, các doanh nghiệp lớn trong nước như CellphoneS, FPT Shop, Thế Giới Di Động đã xây dựng những nền tảng thương mại điện tử hiện đại, đáp ứng đa dạng nhu cầu từ việc tra cứu sản phẩm, so sánh giá, đến đặt hàng trực tuyến. Các hệ thống này được đầu tư bài bản cả về giao diện, trải nghiệm người dùng (UI/UX) lẫn hạ tầng kỹ thuật, mang lại sự tiện lợi và tin cậy cho khách hàng.

Tuy nhiên, khi khảo sát và trải nghiệm thực tế, nhóm nhận thấy rằng các nền tảng thương mại điện tử hiện nay vẫn chủ yếu hoạt động theo cơ chế tương tác truyền thống, nơi người dùng cần thao tác thủ công để tìm kiếm sản phẩm hoặc truy cập thông tin. Một số trang web đã tích hợp chatbot hoặc trợ lý ảo cơ bản, tuy nhiên các hệ thống này chỉ dừng lại ở mức trả lời câu hỏi hoặc hỗ trợ tư vấn đơn giản, chưa có khả năng thực hiện các tác vụ có mức độ tương tác sâu như đặt hàng, hủy đơn hàng, kiểm tra giỏ hàng, hoặc cung cấp thông tin mà khách hàng quan tâm một cách nhanh chóng và an toàn.

Trong bối cảnh trí tuệ nhân tạo (AI) và xử lý ngôn ngữ tự nhiên (Natural Language Processing – NLP) ngày càng phát triển, việc tích hợp Trợ lý Ảo thông minh vào hệ thống thương mại điện tử mở ra hướng tiếp cận mới – nơi người dùng có thể tương tác trực tiếp với hệ thống bằng ngôn ngữ tự nhiên, thay vì qua các bước điều hướng phức tạp. Điều này không chỉ giúp tăng tính cá nhân hóa trong trải nghiệm mua sắm, mà còn giảm tải cho nhân viên chăm sóc khách hàng, đồng thời nâng cao hiệu quả kinh doanh cho doanh nghiệp.

Xuất phát từ thực tế đó, nhóm quyết định lựa chọn và thực hiện đề tài “Hệ thống Thương mại Điện tử cho Bán lẻ Điện thoại tích hợp Trợ lý Ảo AI Agent” với mục tiêu xây dựng một mô hình thương mại điện tử thông minh, có khả năng tương tác tự nhiên, hiểu ngữ cảnh, hỗ trợ và thực hiện hành động thay cho người dùng. Đề tài không chỉ mang ý nghĩa thực tiễn trong việc ứng dụng công nghệ AI vào thương mại điện tử, mà còn góp phần nghiên cứu và phát triển các mô hình tương tác người – máy thế hệ mới.
